\section{Marco de seguridad: RSP y niveles ASL}

\subsection{Dos riesgos centrales}

Dario divide los riesgos de seguridad de la IA en dos categorías principales:

\textbf{Uso indebido catastrófico (QBRN)}: ataques cibernéticos, armas biológicas, armas radiológicas, armas nucleares; capaces de dañar a miles o incluso millones de personas.

\textbf{Riesgo de autonomía}: el modelo actúa por cuenta propia después de obtener más agencia. Es una amenaza más lejana en el tiempo, pero potencialmente más grave.

\begin{warningbox}{La IA rompe el bajo traslape entre «inteligente» y «malvado»}
\textit{«La humanidad siempre ha estado protegida porque el traslape entre las personas verdaderamente inteligentes y bien educadas y las personas que quieren hacer cosas terribles siempre ha sido pequeño.»} La IA puede romper esa correlación: puede convertirse en un «agente más inteligente» que combine un alto coeficiente intelectual con propósitos maliciosos.

Este no es un riesgo del presente, pero \textit{«se acerca a nosotros a una velocidad enorme, como un fantasma.»}
\end{warningbox}

\subsection{El sistema de niveles ASL en detalle}

La Responsible Scaling Policy (RSP) de Anthropic sigue una estructura de «si-entonces»: cuando las pruebas muestran que un modelo posee capacidades peligrosas, se activan los requisitos de seguridad correspondientes.

\begin{table}[H]
\centering
\begin{tabular}{lp{9cm}}
\toprule
\textbf{Nivel} & \textbf{Descripción} \\
\midrule
ASL-1 & Sistemas evidentemente sin riesgo (como el programa de ajedrez Deep Blue) \\
ASL-2 & Sistemas de IA actuales: no son capaces de autorreplicarse y la ayuda que ofrecen para QBRN no supera la de un buscador de Google \\
ASL-3 & Modelos que potencian las capacidades de actores no estatales; requieren medidas de seguridad reforzadas y filtros de despliegue \\
ASL-4 & Modelos que potencian las capacidades de actores estatales o aceleran de forma notable la investigación en IA; pueden requerir interpretabilidad para detectar el «sandbagging» y el engaño \\
ASL-5 & Modelos que superan a toda la humanidad en todas las capacidades de tareas peligrosas \\
\bottomrule
\end{tabular}
\caption{El sistema de niveles AI Safety Levels de Anthropic}
\end{table}

\begin{importantbox}{La filosofía de diseño del ASL}
La RSP fija umbrales de margen para evitar «perder la ventana peligrosa». \textit{«El cuento del lobo es peligroso»}: el riesgo todavía no está aquí hoy, pero se aproxima a una velocidad enorme.

Principio de diseño clave: \textit{«restringir con severidad en cuanto se pueda demostrar que el modelo es peligroso»}, en lugar de imponer límites generales de antemano.
\end{importantbox}

Predicción sobre el cronograma del ASL-3: \textit{«si activamos el ASL-3 en 2025, no me sorprendería en absoluto.»} Para cuando llegue el ASL-4, los modelos podrían ser lo bastante inteligentes para «hacer trampa»: podrían rendir mal de forma deliberada en las pruebas de seguridad (sandbagging) o intentar engañar a quienes los evalúan. Por eso el ASL-4 podría requerir interpretabilidad mecanicista: verificar la seguridad observando directamente el interior del modelo, en lugar de depender de sus propias declaraciones.

\textit{«No definir los criterios concretos del ASL-4 antes de activar el ASL-3»}: esta estrategia ha resultado sensata, porque los desafíos específicos de cada etapa solo pueden entenderse realmente al llegar a esa etapa.

\subsection{Engaño e ingeniería social}

Los modelos de nivel ASL-4 podrían ser lo bastante inteligentes para hacer «sandbagging» de forma deliberada en las pruebas o para engañar a quienes los evalúan. Se necesitan medios de verificación independientes de las propias declaraciones del modelo: ahí es donde entra la interpretabilidad.

\begin{warningbox}{El modelo como «demagogo»}
Una amenaza más sutil es la ingeniería social: el modelo podría volverse extremadamente persuasivo e influir en los ingenieros dentro de la propia empresa. \textit{«En la vida hemos visto muchos ejemplos de demagogos humanos.»} Cuando un modelo es más hábil para persuadir que la mayoría de las personas, los procesos habituales de revisión humana pueden dejar de ser confiables.

Esto apunta a un principio clave: la seguridad no puede depender solo de «preguntarle al modelo si es seguro»; se necesita un mecanismo de verificación independiente, basado en el estado interno.
\end{warningbox}

\subsection{Computer Use: Claude como agente}

Interacción basada en capturas de pantalla: el modelo ve una captura de pantalla y produce coordenadas de clic y acciones de teclado. El ciclo se repite: captura de pantalla $\to$ decisión del modelo $\to$ ejecución $\to$ siguiente captura de pantalla.

\begin{knowledgebox}{«A medio camino de cualquier parte»}
\textit{«Si tienes un modelo preentrenado potente, creo que ya estás a medio camino de cualquier parte.»} El entrenamiento de Computer Use en realidad no fue extenso: es una generalización natural de una capacidad de preentrenamiento potente. La limitación actual son los clics y errores ocasionales, pero Anthropic optó por lanzarlo primero como API (en lugar de dirigirlo directamente al consumidor); Replit fue la empresa que lo desplegó más rápido.
\end{knowledgebox}

En cuanto a las consideraciones de seguridad: Computer Use no es una capacidad fundamentalmente nueva, pero «abre el diafragma» de las capacidades existentes. La inyección de prompts a través del contenido de la pantalla es una superficie de ataque completamente nueva. En el nivel ASL-4, cualquier sandbox podría vulnerarse: \textit{«contener a un modelo malo es muchísimo peor que tener un buen modelo.»}

\subsection{Regulación de la IA}

Anthropic apoya la ley SB 1047 de California (tras sus enmiendas) y es la única empresa de IA que ha expresado una postura detallada y favorable; Google, OpenAI, Meta y Microsoft se oponen todas.

Dario considera que la regulación es necesaria por dos razones:
\begin{enumerate}
    \item Algunas empresas simplemente no cuentan con un mecanismo similar a la RSP: esto es una «externalidad negativa»
    \item No se puede confiar en que las empresas cumplan por voluntad propia con programas voluntarios: \textit{«como industria, nadie nos está vigilando»}
\end{enumerate}

\begin{warningbox}{La paradoja de la regulación}
\textit{«El mayor enemigo de quienes de verdad quieren rendición de cuentas es una regulación mal diseñada.»} Una mala regulación genera un rechazo que produce un consenso antirregulatorio duradero. La regulación debe ser «de precisión quirúrgica, dirigida a los riesgos graves», no una carga general.

\textit{«Si para fines de 2025 seguimos sin haber tomado ninguna medida, eso me preocuparía.»}
\end{warningbox}

\subsection{Resumen de la sección}

La estructura de «si-entonces» de la RSP ofrece un marco progresivo para la inversión en seguridad. El ASL va del 1 al 5, y cada nivel corresponde a requisitos de seguridad distintos. El ASL-3 podría activarse en 2025, y el ASL-4 requerirá interpretabilidad para verificar que el modelo no esté engañando. La regulación debe ser precisa en lugar de general, pero «no hacer nada» es igual de peligroso.
